% GENERATED FILE. Edit canonical lessons, then run npm run generate:books.
% canonical-source-hash: fnv1a64:37e8959a21f1d816
% canonical-lessons: JA-C169-tokidoki, JA-C169-tamani, JA-C169-mettani, JA-C169-maitsuki, JA-C169-kurikaesu, JA-R169-review-164-169

\chapter{Sometimes, Rarely and Every Month}
\label{ch:ja-sometimes-rarely-and-every-month}
\hlchaptermodality{169}

\begin{chapteropening}
\textbf{By the end of this chapter:} \emph{I can say how often I will do things.}

Complete the last lesson of chapter 169: I can say how often I will do things.
\end{chapteropening}

\section[\texorpdfstring{\ja{ときどき} (tokidoki)}{ときどき (tokidoki)}]{\ja{ときどき} (tokidoki) — sometimes}

\label{lesson:JA-C169-tokidoki}



\begin{quote}
Before the new one: say the Japanese for to put off, then the Japanese for to be put off.
\end{quote}

\subsection*{You'll want to know: \ja{ときどき}}
\textbf{\ja{ときどき}} — \emph{tokidoki} — \textquotedblleft{}sometimes\textquotedblright{}.

\begin{quote}
\textbf{\ja{らいねんは} \ja{ときどき} \ja{うみに} \ja{いきます。}} — \emph{rainen wa tokidoki umi ni ikimasu} — \textquotedblleft{}Next year I will sometimes go to the sea.\textquotedblright{}
\end{quote}

\begin{cousinweb}
\textbf{\ja{ときどき}} (\emph{tokidoki}) is native Japanese, built from \emph{toki}, \textquotedblleft{}time\textquotedblright{}, said twice with the second one voiced to \emph{doki}: time and again. It is inherited rather than borrowed (JMdict, KANJIDIC2).
\end{cousinweb}

\begin{grammarlens}[title={What you've built}]
One more word for everyday life at A2.
\end{grammarlens}

\subsection*{Your turn}
\begin{itemize}
\raggedright
  \item \emph{Say it:} \emph{tokidoki}
  \item \emph{Say it:} \emph{tokidoki}, once more
  \item \emph{Say it:} \emph{rainen wa tokidoki umi ni ikimasu} — Next year I will sometimes go to the sea.
  \item \emph{Recall:} say the Japanese for to put off, then the Japanese for to be put off, then say \emph{tokidoki} again
\end{itemize}

\begin{tcolorbox}[breakable,colback=teal!4,colframe=teal!35!black,title={Before you move on}]
What does \ja{ときどき} mean? (\textquotedblleft{}Sometimes\textquotedblright{}.) Say it once more.
\end{tcolorbox}
\section[\texorpdfstring{\ja{たまに} (tamani)}{たまに (tamani)}]{\ja{たまに} (tamani) — once in a while; occasionally}

\label{lesson:JA-C169-tamani}



\begin{quote}
Before the new one: say the Japanese for to be put off, then the Japanese for sometimes.
\end{quote}

\subsection*{You'll want to know: \ja{たまに}}
\textbf{\ja{たまに}} — \emph{tamani} — \textquotedblleft{}once in a while; occasionally\textquotedblright{}.

\begin{quote}
\textbf{\ja{らいげつは} \ja{たまに} \ja{ともだちと} \ja{あいます。}} — \emph{raigetsu wa tamani tomodachi to aimasu} — \textquotedblleft{}Next month I will meet friends once in a while.\textquotedblright{}
\end{quote}

\begin{cousinweb}
\textbf{\ja{たまに}} (\emph{tamani}) is a native Japanese word, inherited rather than borrowed from Chinese, and its older origin is not settled. It is usually written in kana (JMdict).
\end{cousinweb}

\begin{grammarlens}[title={What you've built}]
One more word for everyday life at A2.
\end{grammarlens}

\subsection*{Your turn}
\begin{itemize}
\raggedright
  \item \emph{Say it:} \emph{tamani}
  \item \emph{Say it:} \emph{tamani}, once more
  \item \emph{Say it:} \emph{raigetsu wa tamani tomodachi to aimasu} — Next month I will meet friends once in a while.
  \item \emph{Recall:} say the Japanese for to be put off, then the Japanese for sometimes, then say \emph{tamani} again
\end{itemize}

\begin{tcolorbox}[breakable,colback=teal!4,colframe=teal!35!black,title={Before you move on}]
What does \ja{たまに} mean? (\textquotedblleft{}Once in a while; occasionally\textquotedblright{}.) Say it once more.
\end{tcolorbox}
\section[\texorpdfstring{\ja{めったに} (mettani)}{めったに (mettani)}]{\ja{めったに} (mettani) — rarely (with a negative verb)}

\label{lesson:JA-C169-mettani}



\begin{quote}
Before the new one: say the Japanese for sometimes, then the Japanese for once in a while.
\end{quote}

\subsection*{You'll want to know: \ja{めったに}}
\textbf{\ja{めったに}} — \emph{mettani} — \textquotedblleft{}rarely (with a negative verb)\textquotedblright{}.

\begin{quote}
\textbf{\ja{ちちは} \ja{めったに} \ja{やすみません。}} — \emph{chichi wa mettani yasumimasen} — \textquotedblleft{}My father rarely takes a day off.\textquotedblright{}
\end{quote}

\begin{cousinweb}
\textbf{\ja{めったに}} (\emph{mettani}) is built from \emph{metta} and the native particle \emph{ni}. \emph{Metta} is written with two characters read with their on'yomi, \emph{metsu}, \textquotedblleft{}destroy\textquotedblright{} (Middle Chinese \emph{mjiet}), and \emph{ta}, \textquotedblleft{}many\textquotedblright{} (\emph{ta}); whether it was borrowed from Chinese or the characters were chosen for their sound is not settled (JMdict, KANJIDIC2; Middle Chinese in Baxter's transcription).
\end{cousinweb}

\begin{grammarlens}[title={What you've built}]
One more word for everyday life at A2.
\end{grammarlens}

\subsection*{Your turn}
\begin{itemize}
\raggedright
  \item \emph{Say it:} \emph{mettani}
  \item \emph{Say it:} \emph{mettani}, once more
  \item \emph{Say it:} \emph{chichi wa mettani yasumimasen} — My father rarely takes a day off.
  \item \emph{Recall:} say the Japanese for sometimes, then the Japanese for once in a while, then say \emph{mettani} again
\end{itemize}

\begin{tcolorbox}[breakable,colback=teal!4,colframe=teal!35!black,title={Before you move on}]
What does \ja{めったに} mean? (\textquotedblleft{}Rarely (with a negative verb)\textquotedblright{}.) Say it once more.
\end{tcolorbox}
\section[\texorpdfstring{\ja{まいつき} (maitsuki)}{まいつき (maitsuki)}]{\ja{まいつき} (maitsuki) — every month}

\label{lesson:JA-C169-maitsuki}



\begin{quote}
Before the new one: say the Japanese for once in a while, then the Japanese for rarely.
\end{quote}

\subsection*{You'll want to know: \ja{まいつき}}
\textbf{\ja{まいつき}} — \emph{maitsuki} — \textquotedblleft{}every month\textquotedblright{}.

\begin{quote}
\textbf{\ja{らいねんは} \ja{まいつき} \ja{ははに} \ja{あいます。}} — \emph{rainen wa maitsuki haha ni aimasu} — \textquotedblleft{}Next year I will see my mother every month.\textquotedblright{}
\end{quote}

\begin{cousinweb}
\textbf{\ja{まいつき}} (\emph{maitsuki}) is a mixed word, built from Sino-Japanese \emph{mai}, \textquotedblleft{}every\textquotedblright{}, borrowed from Chinese (Middle Chinese \emph{mwojX}), and native \emph{tsuki}, \textquotedblleft{}moon, month\textquotedblright{} (JMdict, KANJIDIC2; Middle Chinese in Baxter's transcription).
\end{cousinweb}

\begin{grammarlens}[title={What you've built}]
One more word for everyday life at A2.
\end{grammarlens}

\subsection*{Your turn}
\begin{itemize}
\raggedright
  \item \emph{Say it:} \emph{maitsuki}
  \item \emph{Say it:} \emph{maitsuki}, once more
  \item \emph{Say it:} \emph{rainen wa maitsuki haha ni aimasu} — Next year I will see my mother every month.
  \item \emph{Recall:} say the Japanese for once in a while, then the Japanese for rarely, then say \emph{maitsuki} again
\end{itemize}

\begin{tcolorbox}[breakable,colback=teal!4,colframe=teal!35!black,title={Before you move on}]
What does \ja{まいつき} mean? (\textquotedblleft{}Every month\textquotedblright{}.) Say it once more.
\end{tcolorbox}
\section[\texorpdfstring{\ja{くりかえす} (kurikaesu)}{くりかえす (kurikaesu)}]{\ja{くりかえす} (kurikaesu) — to repeat; to do again}

\label{lesson:JA-C169-kurikaesu}



\begin{quote}
Before the new one: say the Japanese for rarely, then the Japanese for every month.
\end{quote}

\subsection*{You'll want to know: \ja{くりかえす}}
\textbf{\ja{くりかえす}} — \emph{kurikaesu} — \textquotedblleft{}to repeat; to do again\textquotedblright{}.

\begin{quote}
\textbf{\ja{らいしゅう} \ja{もういちど} \ja{くりかえします。}} — \emph{raishū mō ichido kurikaeshimasu} — \textquotedblleft{}Next week I will do it over again.\textquotedblright{}
\end{quote}

Polite \textbf{\ja{くりかえします}}, \emph{kurikaeshimasu}; a godan verb.

\begin{cousinweb}
\textbf{\ja{くりかえす}} (\emph{kurikaesu}) is native Japanese, built from \emph{kuri}, \textquotedblleft{}reeling, winding\textquotedblright{}, and \emph{kaesu}, \textquotedblleft{}to send back\textquotedblright{}, both inherited rather than borrowed (JMdict, KANJIDIC2).
\end{cousinweb}

\begin{grammarlens}[title={What you've built}]
One more word for everyday life at A2.
\end{grammarlens}

\subsection*{Your turn}
\begin{itemize}
\raggedright
  \item \emph{Say it:} \emph{kurikaesu}
  \item \emph{Say it:} \emph{kurikaesu}, once more
  \item \emph{Say it:} \emph{raishū mō ichido kurikaeshimasu} — Next week I will do it over again.
  \item \emph{Recall:} say the Japanese for rarely, then the Japanese for every month, then say \emph{kurikaesu} again
\end{itemize}

\begin{tcolorbox}[breakable,colback=teal!4,colframe=teal!35!black,title={Before you move on}]
What does \ja{くりかえす} mean? (\textquotedblleft{}To repeat; to do again\textquotedblright{}.) Say it once more.
\end{tcolorbox}
\section[(dialogue)]{Words from chapters 164–169 — a review}

\label{lesson:JA-R169-review-164-169}



\begin{quote}
Say the Japanese for Tuesday and Wednesday.
\end{quote}

\subsection*{Your turn}
Say these aloud:

\begin{itemize}
\raggedright
  \item Tuesday, Wednesday, Thursday, Friday, Sunday
  \item the seventh of the month, the eighth of the month, the tenth of the month, the twentieth of the month, a birthday
  \item a.m., noon, night-time, this evening, every night
  \item to begin, to start, to finish, to become different, calling off
  \item to be decided, to decide, a deadline, to put off, to be put off
  \item sometimes, once in a while, rarely, every month, to repeat
\end{itemize}

\begin{tcolorbox}[breakable,colback=teal!4,colframe=teal!35!black,title={Before you move on}]
Tuesday? (\textbf{\ja{かようび}}, \emph{kayōbi}.) Wednesday? (\textbf{\ja{すいようび}}, \emph{suiyōbi}.) Thursday? (\textbf{\ja{もくようび}}, \emph{mokuyōbi}.) Friday? (\textbf{\ja{きんようび}}, \emph{kin'yōbi}.) Sunday? (\textbf{\ja{にちようび}}, \emph{nichiyōbi}.) The seventh of the month? (\textbf{\ja{なのか}}, \emph{nanoka}.) The eighth of the month? (\textbf{\ja{ようか}}, \emph{yōka}.) The tenth of the month? (\textbf{\ja{とおか}}, \emph{tōka}.) The twentieth of the month? (\textbf{\ja{はつか}}, \emph{hatsuka}.) A birthday? (\textbf{\ja{たんじょうび}}, \emph{tanjōbi}.) A.m.? (\textbf{\ja{ごぜん}}, \emph{gozen}.) Noon? (\textbf{\ja{ひる}}, \emph{hiru}.) Night-time? (\textbf{\ja{ばん}}, \emph{ban}.) This evening? (\textbf{\ja{こんばん}}, \emph{konban}.) Every night? (\textbf{\ja{まいばん}}, \emph{maiban}.) To begin? (\textbf{\ja{はじまる}}, \emph{hajimaru}.) To start? (\textbf{\ja{はじめる}}, \emph{hajimeru}.) To finish? (\textbf{\ja{おえる}}, \emph{oeru}.) To become different? (\textbf{\ja{かわる}}, \emph{kawaru}.) Calling off? (\textbf{\ja{ちゅうし}}, \emph{chūshi}.) To be decided? (\textbf{\ja{きまる}}, \emph{kimaru}.) To decide? (\textbf{\ja{きめる}}, \emph{kimeru}.) A deadline? (\textbf{\ja{きげん}}, \emph{kigen}.) To put off? (\textbf{\ja{のばす}}, \emph{nobasu}.) To be put off? (\textbf{\ja{のびる}}, \emph{nobiru}.) Sometimes? (\textbf{\ja{ときどき}}, \emph{tokidoki}.) Once in a while? (\textbf{\ja{たまに}}, \emph{tamani}.) Rarely? (\textbf{\ja{めったに}}, \emph{mettani}.) Every month? (\textbf{\ja{まいつき}}, \emph{maitsuki}.) To repeat? (\textbf{\ja{くりかえす}}, \emph{kurikaesu}.)
\end{tcolorbox}
